\chapter{Conclusion}

\section{Key Findings}

Both research efforts reach the same structural conclusion from different
starting points: a budget-constrained Tunisian team cannot win by copying an
international AMR and selling it cheaper, because the components stay imported and
the cost structure cannot be beaten. The defensible advantages are modular
architecture, local manufacturing and support, proximity to the EU, and a
pricing model designed for the local market. Both efforts also agree that
sensor/inspection is the lowest-friction way to win first references, because it
requires no change to a customer's facility.

The two efforts differ mainly on which sector to pair with inspection: the first
emphasises manufacturing intralogistics for Tunisia's automotive exporters; the
second adds universities as an equally strong first market. The merged view below
takes both as the entry wedge and treats manufacturing as the funded mid-term
expansion --- which is consistent with the first effort's own logic of using a
low-friction product to finance the longer manufacturing sales cycle.

\section{Recommendation}

\begin{itemize}
    \item \textbf{First market (now).} Simplified industrial inspection plus
          university pilots, based in Tunisia. These are the two sectors where
          the economic case works without matching an international price ---
          inspection against avoided-incident cost, universities against the cost
          of imported research robots. Both have fast decision cycles and are
          buildable with the team's current ROS2/Nav2 stack, simulation, modular
          chassis, and dashboard.
    \item \textbf{Mid-term (Years 2--3).} Export-oriented automotive and
          electronics manufacturing intralogistics (arm + base) and 3PL/internal
          logistics (base + compartment), sold on a service basis to bypass
          capex resistance.
    \item \textbf{Later.} GCC export of the compartment and sensor
          configurations, where automation budgets are far larger and the modules
          travel without certification risk; private healthcare as a referenceable
          niche; and hospitality only through international chains already present
          in Tunisia.
    \item \textbf{Defer} the shelf-access module to a later, partner-funded phase.
          \textbf{Drop} standalone security and in-store retail as entry products.
\end{itemize}

\section{Cross-Cutting Conditions}

Independent of sector, four conditions apply throughout:
service-based pricing (RaaS, usage-based, and shared deployments between
neighbouring SMEs) \cite{builtin_raas}; a hybrid build that imports only the
irreducible high-value components and fabricates structure and modules locally;
French and Arabic interfaces with a visible local service team; and use of the
Tunisian Startup Act and EU-co-funded industrial-modernization pilots as
financing channels.

\section{Next Steps}

The single most valuable near-term action is to close one credible reference
deployment in inspection or at a university, since both sectors' buyers
benchmark each other. A small number of figures used in this report ---
notably the TAA automotive export and employment figures, the Startup Act
provisions, and the Enova precedent --- still need a citable source attached by
the team before the report is treated as final.
